\section*{Введение}

Современные информационные системы сталкиваются с постоянно растущим числом и сложностью кибератак. Традиционные средства обнаружения, основанные на сигнатурах или статистических аномалиях, не всегда способны эффективно выявлять целенаправленные, многоэтапные атаки, такие как APT. Для повышения точности детекции требуется метод, способный учитывать как структуру системы, так и семантические связи между событиями.

В последние годы широкое распространение получил фреймворк MITRE ATT\&CK~\cite{mitre}, описывающий тактики и техники злоумышленников. Однако прямое сопоставление событий с техниками ATT\&CK затруднено из-за разреженности данных и неоднозначности логов. В этой связи перспективным представляется использование графовых моделей, в частности — атрибутивных метаграфов~\cite{latypov2020,astanin2012}, позволяющих совместно моделировать:

\begin{enumerate}
\item структурные зависимости (узлы = компоненты ИС: процессы, пользователи, файлы);
    \item атрибутивные характеристики (уровень привилегий, сетевые адреса, хэши);
    \item причинно-следственные связи (например, «процесс~A создал процесс~B»).
\end{enumerate}

Атрибутивный метаграф отличается от классического графа возможностью описывать как объекты, так и отношения между ними с использованием атрибутов, что делает его особенно подходящим для моделирования сложных сценариев компрометации. Подход, описанный в настоящей работе, основан на формализации тактик ATT\&CK в виде шаблонов метаграфов и последующем сопоставлении этих шаблонов с динамически строящимся метаграфом системы.

Целью настоящей работы является разработка и экспериментальная оценка подхода к выявлению вредоносной активности на основе атрибутивных метаграфов с интеграцией MITRE ATT\&CK и \seqsplit{совместимостью} с \seqsplit{SIEM-инфраструктурой.}